% ============================================================
%  DSF-AI Layer 6: Topological Constraint Layer (L6-TCL)
%  Deterministic Specification — Final
%  Overleaf-compatible LaTeX
% ============================================================
\documentclass[11pt, a4paper]{article}

\usepackage[margin=2.2cm]{geometry}
\usepackage{amsmath, amssymb, amsthm}
\usepackage{mathtools}
\usepackage{booktabs}
\usepackage{array}
\usepackage{longtable}
\usepackage{xcolor}
\usepackage{hyperref}
\usepackage{titlesec}
\usepackage{enumitem}
\usepackage{fancyhdr}
\usepackage{caption}

\hypersetup{
  colorlinks=true,
  linkcolor=black,
  urlcolor=blue,
}

\definecolor{tfe-gold}{RGB}{180, 140, 60}
\definecolor{tfe-dark}{RGB}{30, 30, 30}
\definecolor{tfe-green}{RGB}{40, 120, 60}
\definecolor{tfe-red}{RGB}{160, 40, 40}
\definecolor{tfe-gray}{RGB}{100, 100, 100}

\pagestyle{fancy}
\fancyhf{}
\rhead{\textcolor{tfe-gray}{\small DSF-AI L6 Topological Constraint Layer}}
\lhead{\textcolor{tfe-gray}{\small Confidential --- Internal Use Only}}
\rfoot{\textcolor{tfe-gray}{\small Page \thepage}}

\titleformat{\section}{\large\bfseries\color{tfe-dark}}{}{0em}{}[\titlerule]
\titleformat{\subsection}{\normalsize\bfseries\color{tfe-dark}}{}{0em}{}
\titleformat{\subsubsection}{\normalsize\itshape\color{tfe-dark}}{}{0em}{}

\setlength{\parindent}{0pt}
\setlength{\parskip}{6pt}

\newtheorem{definition}{Definition}[section]
\newtheorem{axiom}{Axiom}[section]
\newtheorem{theorem}{Theorem}[section]
\newtheorem{corollary}{Corollary}[section]

% ============================================================
\begin{document}

\begin{center}
  {\LARGE \textbf{DSF-AI Layer 6}}\\[6pt]
  {\large Topological Constraint Layer (L6-TCL)}\\[4pt]
  {\large Deterministic Specification --- Final}\\[8pt]
  {\color{tfe-gray}\small Version 1.0 \quad|\quad April 2026 \quad|\quad Confidential}\\[4pt]
  {\color{tfe-gray}\small Joseph Forrester \quad|\quad DSF-AI Architecture}
\end{center}

\vspace{1em}
\hrule
\vspace{1em}

% ============================================================
\section{Introduction}

The Topological Constraint Layer (L6-TCL) is the sixth and highest layer of
the DSF-AI (Deterministic Structural Field AI) architecture. It sits above the
L5 Domain Governance Layer and serves a fundamentally different purpose than
any layer below it.

\textbf{Layers L0--L4} perceive structural state from raw temporal signals.
\textbf{Layer L5} interprets structural state into domain-specific decisions.
\textbf{Layer L6} identifies when structural geometry has collapsed to the
point where only one outcome remains possible --- \textit{Structural Inevitability}.

L6 does not predict. It does not estimate probability. It detects the geometric
condition under which the manifold of possible futures has been reduced to a
singularity. When L6 fires, the outcome is not likely --- it is a mathematical
requirement of the field.

\subsection{Design Constraints}

\begin{enumerate}[label=\arabic*.]
  \item \textbf{Zero heuristics.} Every threshold, constant, and operator in L6
        is derived from the geometry of $n$-dimensional space. No parameter is
        fitted to empirical outcomes.
  \item \textbf{Zero randomness.} No stochastic variables, no random number
        generation, no probabilistic models. All operations are deterministic.
  \item \textbf{Domain agnostic.} L6 operates on abstract dimensional manifolds.
        The domain (financial, medical, autonomous, sensory) is encoded only in
        the Environment Class weighting, not in the core mathematics.
  \item \textbf{Derived from natural constants.} The capture threshold $\tau$ is
        a consequence of the Euler--Gamma relationship in $n$-dimensional
        volume geometry. It is not a design parameter.
\end{enumerate}

% ============================================================
\section{Mathematical Foundations}

This section defines the mathematical objects and operations used throughout
the L6 specification. All notation is self-contained.

\subsection{The $n$-Sphere Volume}

\begin{definition}[$n$-Sphere Volume]
The volume of a unit $n$-sphere of radius $R$ in $n$-dimensional Euclidean space
is given by:
\begin{equation}
  V_n(R) = \frac{\pi^{n/2}}{\Gamma\!\left(\frac{n}{2} + 1\right)} \, R^n
  \label{eq:nsphere}
\end{equation}
where $\Gamma(\cdot)$ is the Euler gamma function.
\end{definition}

\textbf{Key property:} As dimensionality $n$ increases, the volume of the
$n$-sphere concentrates on its shell. The interior (``choice volume'')
progressively empties. This is the mathematical basis for structural
inevitability --- high-dimensional confinement eliminates interior degrees
of freedom.

\subsection{The Gamma Function}

\begin{definition}[Gamma Function]
The Euler gamma function extends the factorial to the real line:
\begin{equation}
  \Gamma(z) = \int_0^{\infty} t^{z-1} e^{-t}\, dt, \qquad z > 0
\end{equation}
For positive integers: $\Gamma(n+1) = n!$
\end{definition}

The gamma function governs how $n$-sphere volume behaves as dimensionality
changes. Its super-exponential growth in the denominator of
Equation~\ref{eq:nsphere} is the mechanism that causes volume collapse at
high $n$. The \textbf{inflection point} of the ratio
$\Gamma(n_{eff}/2 + 1) / \Gamma(n_{start}/2 + 1)$ determines the
geometric capture threshold.

\subsection{Euler's Number ($e$)}

\begin{definition}[Euler's Number]
\begin{equation}
  e = \lim_{n \to \infty} \left(1 + \frac{1}{n}\right)^n \approx 2.71828
\end{equation}
\end{definition}

$e$ appears naturally in the gamma function's asymptotic behavior (Stirling's
approximation: $\Gamma(n+1) \sim \sqrt{2\pi n}\,(n/e)^n$). The critical
ratio $n_{start}/e$ is not a chosen threshold --- it is the point at which
the gamma function's growth rate transitions from polynomial to
super-exponential, causing the volume ratio to exhibit a sharp ``knee.''

% ============================================================
\section{Core Definitions}

\begin{definition}[Active Manifold ($\mathcal{S}_{active}$)]
The set of currently available degrees of freedom (DoF) in the system's
state space. Its volume $V_n(\mathcal{S}_{active})$ represents the
``choice volume'' --- the space of possible future states.
\end{definition}

\begin{definition}[Background Field ($\Phi_{latent}$)]
The omnipresent state-space --- the total manifold within which the active
manifold operates. It is characterized by a \textbf{non-zero gradient}
($\nabla \Phi \neq 0$), representing the fundamental structural imbalance
that drives all flux. The background field is not a passive vacuum; it is
an \textbf{active metric tensor} that shapes material paths by making
non-inevitable trajectories energetically expensive.
\end{definition}

\begin{definition}[Effective Dimensionality ($n_{effective}$)]
The residual degrees of freedom after environmental constraints have been
applied:
\begin{equation}
  n_{effective} = n_{start} - \sum_{i=1}^{k} \operatorname{rank}(C_i)
  \label{eq:neff}
\end{equation}
where $n_{start}$ is the initial dimensionality of the state space and
$C_i$ are the constraint operators imposed by the environment.
\end{definition}

\begin{definition}[Environment Class ($\epsilon$)]
The domain-specific template that defines how constraints are weighted.
Three universal classes are defined:
\begin{center}
\begin{tabular}{@{}lll@{}}
\toprule
\textbf{Class} & \textbf{Weighting} & \textbf{Domain Examples} \\
\midrule
Class-A (Flow)         & Momentum / Velocity   & Market trends, projectile trajectories, fluid dynamics \\
Class-B (Accumulation) & Density / Mass         & Order books, tumor growth, charge accumulation \\
Class-C (Oscillation)  & Frequency / Amplitude  & Signal cycles, heartbeat, thermal oscillation \\
\bottomrule
\end{tabular}
\end{center}
The environment class determines $n_{start}$ and the rank contribution of
each constraint $C_i$.
\end{definition}

% ============================================================
\section{The Core Axiom: Inevitability ($\Omega$)}

\begin{axiom}[Structural Inevitability]
Inevitability is achieved when the active manifold's volume has collapsed
relative to the background field such that no geometrically viable
alternative path exists, and no unresolved exogenous disruption prevents
structural closure.
\end{axiom}

The inevitability scalar $\Omega(t)$ is defined as:

\begin{equation}
  \boxed{
    \Omega(t) = \left[ 1 - \frac{V_n(\mathcal{S}_{active})}{V_n(\Phi_{latent})} \right]
    \cdot \prod_{j=1}^{m} \left(1 - \delta_j^{ext} \cdot \mathcal{H}_{rec}\right)
  }
  \label{eq:omega}
\end{equation}

The first factor measures \textbf{geometric confinement}: as the active
manifold shrinks relative to the background field, $\Omega$ approaches 1.

The second factor is the \textbf{disruption gate}: any active exogenous
disruption ($\delta_j^{ext} = 1$) that has not yet recovered
($\mathcal{H}_{rec} > 0$) multiplies $\Omega$ toward zero, preventing
false locks during structural discontinuities.

% ============================================================
\section{The Exogenous Disruption Operator ($\delta_j^{ext}$)}

\subsection{Motivation}

The system must distinguish between internal structural evolution (which it
tracks through L0--L4) and external discontinuities that originate outside
the observation manifold. These are not random events --- they are
\textbf{deterministically detectable boundary violations}.

\subsection{Definition}

\begin{definition}[Exogenous Disruption Operator]
\begin{equation}
  \delta_j^{ext} = \begin{cases}
    1 & \text{if } \|\Delta \Phi_j\| > \lambda_{disc} \cdot \sigma_{\Phi} \\
    0 & \text{otherwise}
  \end{cases}
  \label{eq:delta}
\end{equation}
where:
\begin{itemize}[nosep]
  \item $\Delta \Phi_j$ is the observed jump in the background field gradient
        at boundary $j$.
  \item $\sigma_{\Phi}$ is the field's own structural variance, derived from
        L3 resonance output. This is \textbf{not} a statistical estimate ---
        it is the deterministic measure of the field's internal variability
        as computed by the L0--L4 kernel.
  \item $\lambda_{disc}$ is a \textbf{frozen rational constant}. It follows
        the same convention as all DSF-AI kernel constants (e.g.,
        $\beta = 37/64$, motion weight $= 3/5$).
\end{itemize}
\end{definition}

\textbf{Interpretation:} The system does not generate randomness. It
\textit{observes} whether the field gradient jumped harder than the field's
own structural variance predicts. If yes, something external disrupted the
manifold. The operator is binary and deterministic.

\subsection{Hysteresis Recovery ($\mathcal{H}_{rec}$)}

\begin{definition}[Hysteresis Recovery]
$\mathcal{H}_{rec}$ is the Quiescence operator applied to the
post-disruption state. It prevents re-locking until the background field
gradient $\nabla \Phi$ has stabilized and the structural closure has
genuinely reformed.

\begin{equation}
  \mathcal{H}_{rec}(t) = \begin{cases}
    1 & \text{if } t - t_{disruption} < \tau_{recovery} \\
    0 & \text{if the manifold has returned to quiescence}
  \end{cases}
\end{equation}

where $\tau_{recovery}$ is the minimum number of observation periods
required for the field gradient to return within $\sigma_{\Phi}$ bounds.
This is measured, not prescribed.
\end{definition}

% ============================================================
\section{The Geometric Capture Threshold ($\tau$)}

\subsection{Derivation from $n$-Dimensional Geometry}

The capture threshold $\tau$ is \textbf{not a design parameter}. It is a
mathematical consequence of how volume behaves in $n$-dimensional space.

\begin{theorem}[Geometric Capture Threshold]
The capture threshold is the volume ratio at which the remaining
$n$-sphere interior can no longer sustain an alternative geometric path:
\begin{equation}
  \boxed{
    \tau = 1 - \frac{\Gamma\!\left(\frac{n_{start}}{2} + 1\right)}
                    {\Gamma\!\left(\frac{n_{effective}}{2} + 1\right)}
    \cdot \pi^{(n_{effective} - n_{start})/2}
  }
  \label{eq:tau}
\end{equation}
\end{theorem}

\subsection{The Gamma-Function Inflection (The ``Knee'')}

\begin{corollary}[Structural Lock Criterion]
The Structural Lock fires when:
\begin{equation}
  n_{effective} < \frac{n_{start}}{e}
  \label{eq:knee}
\end{equation}
where $e \approx 2.71828$ is Euler's number.
\end{corollary}

\textbf{Why $n_{start}/e$?} The gamma function's growth rate transitions
from polynomial to super-exponential at this point (a consequence of
Stirling's approximation). Below this threshold:
\begin{itemize}[nosep]
  \item The $n$-sphere volume collapses toward zero
  \item The shell concentrates all remaining measure
  \item The interior ``choice volume'' is exhausted
  \item The geometric gap between intent and outcome becomes
        \textbf{constraint-dense}
\end{itemize}

This is not a threshold humans or AI systems chose. It is a threshold the
geometry of $n$-dimensional space imposes. The universe itself defines when
escape is impossible.

\subsection{Validation}

The threshold $n_{start}/e$ can be verified numerically for any starting
dimensionality:

\begin{center}
\begin{tabular}{@{}rrrr@{}}
\toprule
$n_{start}$ & $n_{start}/e$ & $V_n/V_{n_{start}}$ at knee & $\tau$ at knee \\
\midrule
8  & 2.94 & 0.0031 & 0.9969 \\
16 & 5.88 & $9.4 \times 10^{-7}$ & $\approx 1.0$ \\
32 & 11.77 & $< 10^{-15}$ & $\approx 1.0$ \\
\bottomrule
\end{tabular}
\end{center}

At higher starting dimensions, the collapse is more extreme. This is
consistent with the physical intuition: more degrees of freedom create
more opportunities for confinement, and the lock is more absolute when
it occurs.

% ============================================================
\section{The Dimensional Grinder}

\subsection{Operational Mechanics}

The Dimensional Grinder is the mechanism that connects the L0--L4 kernel
output to the L6 inevitability calculation. Each constraint $C_i$ observed
by the kernel reduces $n_{effective}$:

\begin{equation}
  n_{effective} = n_{start} - \sum_{i=1}^{k} \operatorname{rank}(C_i)
\end{equation}

\textbf{Constraint sources from L0--L4:}
\begin{center}
\begin{tabular}{@{}lll@{}}
\toprule
\textbf{Kernel Output} & \textbf{Constraint Type} & \textbf{Rank Contribution} \\
\midrule
$S_{UF}$ (Stability) & Structural convergence & $\operatorname{rank} \propto S_{UF}$ \\
$D_k$ (Displacement) & Directional confinement & $\operatorname{rank} = |D_k|$ (0 or 1) \\
$M_k$ (Motion) & Momentum lock & $\operatorname{rank} \propto |M_k|$ \\
$R_{rev,k}$ (Reversal) & Path elimination & $\operatorname{rank} = R_{rev,k}$ (0 or 1) \\
$C_k$ (Cohesion) & Structural binding & $\operatorname{rank} \propto C_k / (1 + C_k)$ \\
$P_k$ (Pressure) & Environmental compression & $\operatorname{rank} \propto P_k / (1 + P_k)$ \\
$B_k$ (Breathing) & Conviction saturation & $\operatorname{rank} \propto |B_k|$ \\
$U^*_k$ (Uncertainty) & Residual freedom & Negative: increases $n_{eff}$ \\
\bottomrule
\end{tabular}
\end{center}

Note that $U^*_k$ (uncertainty) \textbf{increases} effective dimensionality
--- it represents residual freedom that the field has not yet resolved.
This is the system's self-correcting mechanism: high uncertainty prevents
premature lock.

\subsection{The Logic of Exhaustion}

As constraints accumulate:
\begin{enumerate}[nosep]
  \item $n_{effective}$ decreases
  \item $V_n(n_{effective})$ collapses (volume concentrates on shell)
  \item The ratio $V_n(\mathcal{S}_{active}) / V_n(\Phi_{latent})$ approaches zero
  \item $\Omega$ approaches 1
  \item When $n_{effective} < n_{start}/e$, the Structural Lock fires
\end{enumerate}

This is the ``telegraph'' --- the inference of constraints. The system
does not decide when inevitability occurs. The geometry announces it.

% ============================================================
\section{The Structural Lock (SL-1)}

\subsection{Output Specification}

When $\Omega(t) > \tau$ and $n_{effective} < n_{start}/e$, L6 issues a
\textbf{Structural Lock (SL-1)} to the L5 Domain Governance Layer.

SL-1 is not a signal. It is not a recommendation. It is a \textbf{geometric
certification} that the manifold of possible futures has collapsed to a
singularity. The outcome is no longer probable --- it is structurally
required.

\subsection{Relationship to L5}

\begin{center}
\begin{tabular}{@{}lll@{}}
\toprule
\textbf{Layer} & \textbf{Output} & \textbf{Nature} \\
\midrule
L5 (Domain Governance) & Accumulate / Hold / Avoid & Structural interpretation \\
L6 (Topological Constraint) & Structural Lock (SL-1) & Geometric certainty \\
\bottomrule
\end{tabular}
\end{center}

SL-1 \textbf{overrides} L5 signaling. When the geometry has locked, the
domain governance layer's interpretation is superseded by topological
certainty.

% ============================================================
\section{The Background Field as Active Metric Tensor}

The background field $\Phi_{latent}$ is not a passive vacuum. It is an
\textbf{active metric tensor} that shapes the material manifold by making
non-inevitable paths energetically expensive.

\begin{definition}[Negative Space Potential Well]
The ``Negative Space'' --- the quiescent, structurally compressed region
of the field --- acts as a potential well. The longer the system remains
in quiescence ($\tau_{in}$), the deeper the well, and the higher the
topological pressure on the material manifold.
\end{definition}

This connects directly to the Quiescence operator defined in the L5.5
Advanced Localized Thermodynamics specification:
\begin{itemize}[nosep]
  \item Potential energy: $E_p \propto \tau_{in}$
  \item Exhaustion timer: $\tau_{out} = \lfloor \tau_{in} / 3 \rfloor$
  \item Asymmetric exit: $T_{exit} = P_{apex} - 2\epsilon$
\end{itemize}

In L6 terms, quiescence is the mechanism by which the background field
\textit{actively reduces} $n_{effective}$. Each bar of structural
compression is a constraint $C_i$ applied by the field itself. The
Negative Space is the ``Something Else in the shadows'' --- the
mathematical requirement that makes the material outcome inevitable.

% ============================================================
\section{Domain Applications}

L6-TCL is domain-agnostic. The core mathematics ($\Omega$, $\tau$, the
Dimensional Grinder) are identical across all domains. Only the
Environment Class weighting and constraint mapping change.

\subsection{Finance (TFE --- Current Proving Ground)}

\begin{itemize}[nosep]
  \item Environment Class: \textbf{A (Flow)} --- momentum/velocity weighted
  \item $n_{start}$: 8 (the kernel's structural state vector)
  \item Constraints $C_i$: mapped from $D_k$, $S_{UF}$, $M_k$, $R_{rev,k}$,
        $C_k$, $P_k$, $B_k$, $U^*_k$
  \item SL-1 action: maximum position sizing with full conviction
  \item Interpretation: the structural geometry of the price field has
        locked --- the move is a geometric necessity, not a probabilistic setup
\end{itemize}

\subsection{Oncology}

\begin{itemize}[nosep]
  \item Environment Class: \textbf{B (Accumulation)} --- density/mass weighted
  \item $n_{start}$: determined by the biological signal dimensionality
  \item Constraints $C_i$: mapped from cellular signal structural state
  \item SL-1 action: identification of \textbf{pathological irreversibility}
  \item Interpretation: the geometry of the cellular environment has locked
        into a malignant trajectory --- the tumor is now a mathematical
        requirement of the field
\end{itemize}

\subsection{Autonomous Systems}

\begin{itemize}[nosep]
  \item Environment Class: \textbf{A (Flow)} --- momentum/velocity weighted
  \item $n_{start}$: determined by the environmental sensor array dimensionality
  \item Constraints $C_i$: mapped from spatial-temporal structural state
  \item SL-1 action: identification of \textbf{collision singularity} or
        \textbf{safe passage certainty}
  \item Interpretation: the moment where no steering or braking input can
        alter the outcome because the Negative Space of the environment
        has reached $n \to 1$
\end{itemize}

\subsection{Sensory Cognition (ArcLoom / Krimelacks)}

\begin{itemize}[nosep]
  \item Environment Class: \textbf{C (Oscillation)} --- frequency/amplitude weighted
  \item $n_{start}$: determined by the number of active sensory kernel clusters
  \item Constraints $C_i$: inter-kernel resonance convergence
  \item SL-1 action: \textbf{perceptual certainty} --- the sensory field has
        resolved to a single structural interpretation
  \item Interpretation: the system ``knows'' what it perceives, not because
        it classified the input, but because the geometry of the sensory
        field has collapsed to one possibility
\end{itemize}

% ============================================================
\section{Invariants and Constraints}

\begin{enumerate}
  \item $\Omega(t) \in [0, 1]$. Values outside this range indicate
        implementation error.
  \item $n_{effective} \geq 0$. Negative effective dimensionality is
        undefined and must be clamped.
  \item $\tau$ is \textbf{never empirically tuned}. It is computed from
        Equation~\ref{eq:tau} using the gamma function. Any implementation
        that substitutes a fitted constant is non-compliant.
  \item $\delta_j^{ext}$ is \textbf{never generated randomly}. It is a
        deterministic boundary detection. Any implementation that introduces
        random number generation is a violation.
  \item $\lambda_{disc}$ is a frozen rational constant. Once set, it is
        not modified based on domain performance.
  \item SL-1 is irrevocable within the current observation window. Once
        the Structural Lock fires, L5 cannot override it. Only a new
        exogenous disruption ($\delta_j^{ext} = 1$) can release the lock
        after $\mathcal{H}_{rec}$ clears.
  \item L6 does not modify L0--L4 kernel output. It is a \textbf{read-only}
        consumer of the structural state vector.
\end{enumerate}

% ============================================================
\section{Academic Integrity Statement}

The L6-TCL specification is anchored to two natural mathematical constants:

\begin{enumerate}
  \item \textbf{Euler's number ($e$)} --- determines the dimensional
        inflection point at which volume collapse becomes irreversible.
  \item \textbf{The Gamma function ($\Gamma$)} --- governs the relationship
        between dimensionality and volume, producing the capture threshold
        as a mathematical consequence rather than a design choice.
\end{enumerate}

By deriving $\tau$ from the inflection point of the gamma function, L6
removes all human and AI design decisions from the inevitability
determination. The system does not \textit{decide} when an outcome is
inevitable. It \textit{observes} when the geometry of $n$-dimensional
space has exhausted all alternatives.

\textbf{Zero heuristics. Zero randomness. Zero fitted parameters.}

The threshold is a derivative of $n$-dimensional geometry. The disruption
operator is a deterministic measurement of boundary discontinuities. The
mathematics is scale-invariant --- whether $n$ represents market variables,
cellular signaling pathways, or sensory kernel clusters, the gamma-function
inflection point remains constant.

% ============================================================
\section{Relationship to DSF-AI Architecture}

\begin{center}
\begin{tabular}{@{}llll@{}}
\toprule
\textbf{Layer} & \textbf{Function} & \textbf{Output} & \textbf{Nature} \\
\midrule
L0 & SEV series & Normalized temporal signal & Perception \\
L1 & Boundary operator & Gate segmentation & Structure \\
L2 & Regime interpretation & $\chi/\psi$ classification & Classification \\
L3 & Resonance & Cross-scale coherence & Integration \\
L4 & Decision dynamics & $D_k, M_k, B_k, \ldots$ & State vector \\
L5 & Domain governance & Accumulate / Hold / Avoid & Interpretation \\
\textbf{L6} & \textbf{Topological constraint} & \textbf{Structural Lock (SL-1)} & \textbf{Certainty} \\
\bottomrule
\end{tabular}
\end{center}

Each layer builds on the one below. L6 is the capstone: it takes the
full L0--L4 structural state vector, maps it through the Dimensional
Grinder, and determines whether the geometry has reached the point of
no return. When it has, the physics speaks --- and the outcome is not
a prediction but a requirement.

% ============================================================
\section{Version History}

\begin{center}
\begin{tabular}{@{}lll@{}}
\toprule
\textbf{Version} & \textbf{Date} & \textbf{Description} \\
\midrule
0.1 & 2026-04-09 & Initial draft. Core axiom, disruption operator, environment classes. \\
0.2 & 2026-04-09 & Replaced ``stochastic'' with Exogenous Disruption Operator ($\delta_j^{ext}$). \\
    &            & Derived $\tau$ from gamma-function inflection ($n_{start}/e$). \\
1.0 & 2026-04-09 & Final specification. Full mathematical framework, domain applications, \\
    &            & constraint mapping, academic integrity statement. \\
\bottomrule
\end{tabular}
\end{center}

\vspace{2em}
\hrule
\vspace{0.5em}
{\color{tfe-gray}\small This document is confidential and for internal use only.
It constitutes the deterministic specification for DSF-AI Layer 6: the
Topological Constraint Layer. All mathematics are derived from natural
constants ($e$, $\Gamma$) and $n$-dimensional geometry. No parameters
are fitted, estimated, or heuristically determined. This is the Physics
of Certainty.}

\end{document}
